\section{Complete compiled-policy services and economic comparisons}\label{sec:study48}
The new study separates three questions: the cost of representing a fixed NBO policy, the complete cost of constructing and certifying a policy, and the actual cost of using one controller rather than another. A representation benchmark cannot by itself answer either of the latter questions. Conversely, a conventional comparator receives the same innovation law, feasibility conditions and all-state comparison principle, not a deliberately inefficient implementation chosen to favor the neural operator.

\subsection{Design, stopping rules and conventional baselines}
The design was fixed before the complete execution. The principal catalogue uses the original constrained two-state economy of Section~\ref{sec:constrained47}, horizons $T=2,3$, prices $p=1,4$, and $N=M\in\{4,8,16,32,64\}$. The three methods are the compiled cone witness, uniform tensor FVI and error-driven nonuniform tensor FVI. There are three isolated process repetitions for each method and economic cell: 36 services and 180 rungs. Every rung starts from primitives and uses its own future continuation. All rungs are executed, including those after a target crossing; first-target work charges the complete prefix, including every prior failure.

The conventional adaptive method forms an own-future pilot on a $5\times5$ grid, uses its coordinatewise second differences to allocate nonuniform dyadic knots, and certifies with the actual largest coordinate gaps. Its pilot, allocation and repeated target evaluations are charged. The section on nonuniform fitted-value comparison in the supplement specifies the fixed heuristic and proves the strengthened nearest-node certificate on these unequal grids. It is a conventional adaptive-coordinate FVI comparator, not a sparse-grid method or an unrestricted optimal adaptive partition.

All methods share outward arithmetic, feasible action fractions and midpoint integration with its deterministic remainder under the continuous uniform law. Each service includes a finite acquired-state deployment checksum and durable checkpoint output. Complete prefix wall and CPU times, method-specific evaluation counts, coefficient bit lengths, stored actor scalars, checkpoint bytes and peak resident memory are recorded. The internal clock excludes startup and the common non-economic warm-up; a separate process clock includes them. No independent cost comparison is hidden inside that boundary: its joint service is reported separately below.

\begin{table}[htbp]
\centering
\caption{Complete two-state construction services at the declared cap}\label{tab:core48}
\footnotesize
\begin{tabular}{rrcrrrr}
\toprule
$T$ & $p$ & Method & Loss bound & Median (s) & Range (s) & Peak MiB\\
\midrule
2 & 1 & C & 0.579 & 48.001 & [47.981, 49.481] & 92.0 \\
2 & 1 & U & 0.719 & 36.308 & [36.280, 36.426] & 74.5 \\
2 & 1 & A & 0.719 & 36.689 & [36.685, 36.696] & 74.6 \\
2 & 4 & C & 0.592 & 48.077 & [47.790, 48.600] & 91.3 \\
2 & 4 & U & 0.727 & 36.391 & [36.389, 37.257] & 73.5 \\
2 & 4 & A & 0.727 & 37.025 & [36.669, 37.110] & 73.7 \\
3 & 1 & C & 0.830 & 71.933 & [71.549, 72.464] & 93.1 \\
3 & 1 & U & 1.003 & 54.649 & [54.146, 54.650] & 74.8 \\
3 & 1 & A & 1.003 & 55.677 & [55.256, 57.886] & 75.5 \\
3 & 4 & C & 0.853 & 71.963 & [71.798, 72.082] & 93.5 \\
3 & 4 & U & 1.013 & 54.151 & [53.822, 54.812] & 74.9 \\
3 & 4 & A & 1.013 & 55.026 & [54.930, 55.110] & 75.3 \\
\bottomrule
\end{tabular}
\begin{minipage}{0.97\linewidth}\footnotesize\vspace{3pt}
C denotes compiled witness, U uniform tensor FVI, and A error-driven coordinate FVI. Each entry ends at $N=M=64$ and includes the complete five-rung prefix through durable record output. The range is over three isolated repetitions with identical checkpoints, not three economic draws. Peak resident memory includes the interpreter and warm-up. Bounds are displayed to three decimals; exact rational tests determine attainment.
\end{minipage}
\end{table}

\begin{table}[htbp]
\centering
\caption{Prospective target attainment in the two-state catalogue}\label{tab:attainment48}
\footnotesize
\begin{tabular}{crrrrr}
\toprule
Method & $4$ & $2$ & $1$ & $1/2$ & $1/4$\\
\midrule
C & 12/12 & 12/12 & 12/12 & 0/12 & 0/12 \\
U & 12/12 & 12/12 & 6/12 & 0/12 & 0/12 \\
A & 12/12 & 12/12 & 6/12 & 0/12 & 0/12 \\
\bottomrule
\end{tabular}
\begin{minipage}{0.97\linewidth}\footnotesize\vspace{3pt}
All failures and every prior rung are retained. Counts concern twelve deterministic services per method, crossing four economic cells and three timing repetitions. A repeated failure is not an independent estimate of optimizer reliability. No historical failure is reclassified.
\end{minipage}
\end{table}

Against uniform FVI, compiled witness is earlier in 3 of 30 common successful target/repetition comparisons, later in 27, and tied in 0.
Against error-driven FVI, compiled witness is earlier in 15 of 30 common successful target/repetition comparisons, later in 15, and tied in 0.
These are observed complete-prefix wall times, not a hardware-general speed law. The comparison uses the prospective targets rather than retrospectively selecting a favorable breakpoint. The full positive-tolerance partition, including one-sided and joint nonattainment, is reported in the supplement and machine-readable frontier.
The error-driven rule returns uniform grids in all 210 recorded date models; all 60 adaptive checkpoints have exactly the same continuations and actors as their uniform-FVI counterparts. Thus its additional observed cost is pilot and allocation overhead, not evidence of a realized nonuniform-grid improvement. The nonuniform verification theorem is proved and tested, but an effective spatial-adaptation advantage is not established by this catalogue.
One prospective tolerance illustrates the distinction between common-resolution speed and target work. At horizon three, price four and target two, witness crosses at N=32 while uniform FVI first crosses at N=64. Median complete-prefix times are 5.498 and 54.146 seconds, respectively. This is one economic cell with three timing repetitions. At target one, witness attains all twelve services and each FVI implementation attains six; none attains one half or one quarter at its cap.


\subsection{Representation identity and higher-dimensional execution}
The fixed-object diagnostic uses all twelve distinct original constrained witness checkpoints, crossing $T=2,3$, $p=1,4$ and $N=4,8,16$. Both the original dense evaluator and the compiler evaluate the same fixed point catalogue. Compiler setup is charged. Exact rational off-grid checks verify values and original winning indices, including the deterministic tie convention; outward machine intervals are also checked. These clocks do not include new target generation and are not added to historical training times.
\begin{table}[htbp]
\centering
\caption{Fixed-policy representation costs, including compiler setup}\label{tab:compiler48}
\footnotesize
\begin{tabular}{rrrrrr}
\toprule
$N$ & Observations & Median ratio & Minimum & Maximum & Compiled (s)\\
\midrule
4 & 12 & 1.93 & 1.87 & 2.37 & 0.0074 \\
8 & 12 & 2.92 & 2.80 & 4.10 & 0.0137 \\
16 & 12 & 4.48 & 4.42 & 4.52 & 0.0385 \\
\bottomrule
\end{tabular}
\begin{minipage}{0.97\linewidth}\footnotesize\vspace{3pt}
The ratio is dense load-plus-query time divided by compiled setup-plus-query time, over the same original continuation and points. Each resolution crosses four economic cells and three repetitions. This is a same-policy representation experiment, not a construction-to-certificate speedup or an independently trained native comparator. Exact witness-index checks accompany the numerical interval checks.
\end{minipage}
\end{table}


To expose, rather than suppress, dimension factors, a cycle-coupled extension executes $d=3,4$ with continuous common uncertainty, state-dependent capacity and a nontrivial scalar action. Its normalized stage and terminal costs recover the two-state economy exactly at $d=2$. The full primitive formulas and valid moduli appear in the supplement. The stress catalogue fixes $p=1$, $T=2,3$, and the compiled witness and uniform multilinear FVI methods. Dimension three uses $N=M=4,8,16$; dimension four uses $N=M=2,4,8$. This adds eight services and 24 rungs.
\begin{table}[htbp]
\centering
\caption{Executed nonlinear dimension stress cases}\label{tab:dimension48}
\footnotesize
\begin{tabular}{rrcrrrrr}
\toprule
$d$ & $T$ & Method & $N$ & States & Loss bound & Seconds & MiB\\
\midrule
3 & 2 & C & 16 & 4913 & 2.330 & 9.505 & 64.9 \\
3 & 2 & U & 16 & 4913 & 2.814 & 6.677 & 52.3 \\
3 & 3 & C & 16 & 4913 & 3.349 & 13.906 & 66.5 \\
3 & 3 & U & 16 & 4913 & 3.915 & 9.905 & 52.6 \\
4 & 2 & C & 8 & 6561 & 4.630 & 7.733 & 69.7 \\
4 & 2 & U & 8 & 6561 & 5.456 & 7.368 & 50.3 \\
4 & 3 & C & 8 & 6561 & 6.642 & 11.550 & 71.2 \\
4 & 3 & U & 8 & 6561 & 7.589 & 10.982 & 51.2 \\
\bottomrule
\end{tabular}
\begin{minipage}{0.97\linewidth}\footnotesize\vspace{3pt}
Price is one. Each service includes all three declared rungs and has one repetition. The maximum resolution differs by dimension and is not a matched-accuracy scaling design. The common innovation remains continuous; its quadrature remainder, state-dependent capacity and all-state loss account are retained.
\end{minipage}
\end{table}

The smaller four-dimensional cap is explicit. These are memory-limited stress designs with one observation per cell, not repeated fixed-accuracy scaling estimates. They test an executed nonlinear implementation in dimensions above two without turning the finite-dimensional theorem into a dimension-free performance claim. The retained tensor state count, action factors, innovation work and $2^d$ corner terms determine what this implementation can afford.

\subsection{Direct constrained-policy decisions and information prices}
The direct comparison freezes the original $N=16$ witness and FVI policies. It does not select the new R48 policies after examining their costs. Forty-eight policy-pair tasks cross the four economic cells, four initial laws and three observation contracts of Section~\ref{sec:economic48}. The full study contains $48\times262144=12582912$ paired paths. Numerical path widths, every unresolved actor or sensor branch, endpoint moments, simultaneous intervals and decisions are deposited.
\begin{table}[htbp]
\centering
\caption{Direct witness-minus-FVI constrained-policy costs}\label{tab:direct48}
\footnotesize
\begin{tabular}{rrclrr}
\toprule
$T$ & $p$ & Sensor & Interval hull over four laws & Signed & No fee recovery\\
\midrule
2 & 1 & exact & [-0.001247, 0.002717] & 1/4 & 4/4 \\
2 & 1 & 6 & [-0.001113, 0.002716] & 1/4 & 4/4 \\
2 & 1 & 10 & [-0.001220, 0.002717] & 1/4 & 4/4 \\
2 & 4 & exact & [-0.001142, 0.001840] & 0/4 & 4/4 \\
2 & 4 & 6 & [-0.001014, 0.001602] & 0/4 & 4/4 \\
2 & 4 & 10 & [-0.001151, 0.001832] & 0/4 & 4/4 \\
3 & 1 & exact & [-0.001501, 0.002627] & 1/4 & 4/4 \\
3 & 1 & 6 & [-0.001364, 0.002528] & 1/4 & 4/4 \\
3 & 1 & 10 & [-0.001455, 0.002639] & 1/4 & 4/4 \\
3 & 4 & exact & [-0.001356, 0.001806] & 0/4 & 4/4 \\
3 & 4 & 6 & [-0.001342, 0.001783] & 0/4 & 4/4 \\
3 & 4 & 10 & [-0.001365, 0.001808] & 0/4 & 4/4 \\
\bottomrule
\end{tabular}
\begin{minipage}{0.97\linewidth}\footnotesize\vspace{3pt}
Each hull contains the four separate simultaneous intervals for the declared initial laws; it is not an all-state band. There are 48 contrasts, each with 262,144 common interval paths, and family error at most 0.01 under the independent-bin model. Signed counts identify intervals excluding zero. The final column counts intervals strictly inside $(-1/64,1/64)$, ruling out recouping either replacement fee. Display endpoints are rounded outward. The observation fee cancels when the two controllers use the same sensor.
\end{minipage}
\end{table}

Table~\ref{tab:direct48} summarizes interval hulls across the four declared initial laws; its rows are not all-state confidence bands. The supplement reports every interval separately. The zero threshold and the replacement-fee threshold $1/64$ are different decision questions. A no-recoupment result is conditional on that predeclared theoretical fee. It neither establishes exact policy equality nor measures empirical welfare equivalence.

The sensor-price experiment evaluates all thirteen precisions $b=4,\ldots,16$ at each of the three declared information prices, for every original witness checkpoint. Table~\ref{tab:sensors48} gives the finest original resolution; the evidence includes all coarser checkpoints and every candidate precision, not only the minimizing choice.
\begin{table}[htbp]
\centering
\caption{Priced precision for the original finest witness policies}\label{tab:sensors48}
\footnotesize
\begin{tabular}{rrcrrrr}
\toprule
$T$ & $p$ & Bit price & $b_*$ & Acquisition & Fee & Total account\\
\midrule
2 & 1 & $\frac{1}{16384}$ & 16 & 0.0027 & 0.0038 & 2.3212 \\
2 & 1 & $\frac{1}{4096}$ & 14 & 0.0040 & 0.0133 & 2.3320 \\
2 & 1 & $\frac{1}{1024}$ & 12 & 0.0092 & 0.0455 & 2.3694 \\
2 & 4 & $\frac{1}{16384}$ & 16 & 0.0034 & 0.0038 & 2.3763 \\
2 & 4 & $\frac{1}{4096}$ & 14 & 0.0047 & 0.0133 & 2.3871 \\
2 & 4 & $\frac{1}{1024}$ & 12 & 0.0101 & 0.0455 & 2.4247 \\
3 & 1 & $\frac{1}{16384}$ & 16 & 0.0043 & 0.0056 & 3.3310 \\
3 & 1 & $\frac{1}{4096}$ & 14 & 0.0063 & 0.0193 & 3.3468 \\
3 & 1 & $\frac{1}{1024}$ & 12 & 0.0145 & 0.0661 & 3.4017 \\
3 & 4 & $\frac{1}{16384}$ & 16 & 0.0054 & 0.0056 & 3.4215 \\
3 & 4 & $\frac{1}{4096}$ & 15 & 0.0061 & 0.0207 & 3.4374 \\
3 & 4 & $\frac{1}{1024}$ & 13 & 0.0103 & 0.0716 & 3.4925 \\
\bottomrule
\end{tabular}
\begin{minipage}{0.97\linewidth}\footnotesize\vspace{3pt}
The original $N=16$ checkpoint is fixed. Prices are per coordinate-bit per date; $b$ ranges from four to sixteen and action spacing is $1/4096$. The selected bit count minimizes the exact augmented uniform upper account, not measured policy cost. All thirteen candidate precisions and the coarser original checkpoints remain in the evidence. Displayed nonnegative accounts are rounded upward.
\end{minipage}
\end{table}

The optimum minimizes the uniform augmented upper account. Since the full-information certificate $G$ is unchanged across observation precisions, the integer choice responds to the objective sensitivity and the price of information, not to an arbitrary selection of an attractive state-error radius. A coarse state cover can still make the total account large; paying for more bits cannot repair that separate construction error.

\subsection{What changes in the evidence}
The compiler theorem supplies a constructive work reduction for a defined NBO backend without changing its policy. The complete catalogue measures that backend against uniform and error-driven conventional FVI with all failed resolutions included. Direct constrained costs and priced observation support explicit economic decisions for fixed policies under declared information contracts. These results extend the original program, rather than replace it with a different subject.

The historical adverse findings remain: the original nonlinear optimizer catalogue, selected residual-only recertifications, unresolved neural--ridge signs, and mixed-precision timing records are not reclassified. Identical native and ReLU envelopes remain one mathematical construction. The new deterministic repetitions do not estimate a population success probability for nonconvex optimization; the direct confidence statements concern their fixed policies under the sampling model; the stress clocks do not identify hardware-general scaling. The comparative claims in this article are restricted to the objects actually proved or executed, while the controlled-economy and application-specific theoretical scope remains intact.
